
\thispagestyle{empty}
\vspace*{0.30\textheight}
\begin{quote}\kaishu
  给他：\\
  献给他，那位感受力含蓄内敛的诗人，怀着一股内在的火热锋芒，而首要的，
  是他对创造了这片故土的罗马尼亚灵性的深刻理解。谨向这位跻身
  历代大公之列的伟大爱国者致敬。

  给她：\\
  献给她，她的正义感、感受力与人格威望，造就了足球也同样需要的那种
  框架，一如任何要把人们连同他们的精力与热忱聚拢起来的尝试。
  谨向她致以深深的敬意：她在体育运动中付出努力，迎着一切障碍而上，
  只为理解真相。
\end{quote}
\clearpage

\chapter*{前言}
\addcontentsline{toc}{chapter}{前言}
\markboth{前言}{前言}%
\setcounter{chapter}{0}%

\trFloerHomology{}（Floer homology）是当今\trSymplectic\trTopology{}中的一项
基本技术。它的灵感来自1980年代\trWitten{}与\trGromov{}的思想，已经使
许多困难问题的求解成为可能，而且还在继续这样做。

本书致力于解决其中一个问题，即那条出自\trArnold{}的著名猜想：它主张
用一个\trInvariant{}给出\trHamiltonian{}系统的\trPeriodic\trTrajectory{}数目
的下界，而这个\trInvariant{}只依赖于该系统在其中演化的\trSymplectic\trManifold{}的\trTopology。这个下界与著名的\trMorse{}不等式极为相似——后者给出一个函数的%
\trCriticalPoint{}数目的下界。这种相似绝非偶然：\trFloerHomology{}正是%
\trManifold\trHomology{}的（无穷维）类似物，后者由\trMorseComplex{}按“\trWitten{}式”
的办法计算出来；两种情形里起主要作用的，都是连接诸\trCriticalPoint{}的%
\trTrajectory{}的%
\trModuliSpace{}（对\trMorseHomology{}而言是函数的\trCriticalPoint，
对\trFloerHomology{}而言是泛函的\trCriticalPoint）。

2004--2005年，我们为一年级与二年级的研究生开了一门课，毋宁说是两门
课，讲的正是这些内容。我们自然从\trMorse{}理论讲起。我们俩都非常
喜欢\trMilnor{}的那本书，我们两人都从中领教了\trMorseFunction{}的
存在性、丰富性，尤其是它们何等有用；于是我们开始为学生写讲义，这
相当容易……

随后事情就变得更难了——没有哪本书以更现代的观点讲述\trMorseHomology，
即用\trTrajectory{}空间来定义\trMorseComplex{}、给出其构造并证明其
不变性，从而让我们得以进而构造\trFloerComplex{}。我们再也不能照抄了；
如今得动用一点想象力了。

在把课程的第一部分讲到令学生们满意之后，我们向\trFloerHomology{}进军。
我们这些\trTopology{}学家与几何学家每天都在使用的那些\trMorseHomology{}的对象
与技术，被改造成了\trFloerHomology{}的对象与技术。这个理论的魅力
——或者说是魅力之一——及其力量，在于它除了几何与\trTopology{}之外，还要用到
大量的分析、\trFredholmOperator{}与\trSobolevSpace{}。要向真正的学生
讲清楚这一点并非易事。正因此，我们决定把讲义继续写下去。

尽管许多研究著作已经而且仍在使用这些技术，眼下也有许多学生需要它们，
关于这个主题却始终不曾有过一本大体自足的书。

五年过去了。在这五年里，我们打磨、订正、加长、细化，添入上界、下界、
等式与比较；陈述并证明了七十三条定理、一百二十一条命题、一百零三条
引理，画出九十八幅插图\footnote{见脚注~2。}（还布置了一定数量的
练习——它们一反惯例，不含任何“留给读者作为练习”的重要结果的
证明）……

五年过去了。在这五年里，另一些学生读了这些讲义并提出了意见，
使我们确信：这些讲义满足了一种需要，若不把它们再改进、写成一部书，
就未免太傻了。

书就在这里——一部献给“\trMorse--\trFloer{}式”\trHomology{}方法之力量与荣光
的书。

\noindent\textbf{你需要知道的事……}

假装能写出一本自足的书，实属不易。尽管如此，本书源自一门开给真正的
研究生的课——为了他们，我们首先得“回忆”一下什么是\trManifold{}。
心里惦记着他们和他们的同学，我们把“你需要知道的事”汇总在书末的
三个附录里：即我们所用到的\trDifferential{}几何、代数\trTopology{}与分析的基础结果。其中
有时给出完整的定义和（或）证明，有时只给些提示。索引应该能帮助读者
找到自己的路。

\noindent\textbf{致谢}

感谢所有修过这门课的学生，特别是埃米莉·布尔贡德、奥利维耶·多达纳、Shanna Li、亚历山大·穆通、Emmanuel Rey、Nelson Souza。

感谢阿涅丝·加德布莱，她细致地通读了本书的许多预备版本。感谢
克莱芒丝·拉布鲁斯与樊尚·于米利耶尔，感谢他们的问题与建议。
感谢 Andr\'e Carneiro，感谢他提示的种种修正。感谢埃马纽埃尔·奥普什泰因，
感谢他帮我们找到的那些答案。

感谢弗朗索瓦·\trLaudenbach，杜萨·\trMcDuff。感谢让-克洛德·\trSikorav，
感谢他贴近而热忱的通读、写满数页的意见，以及我们与他之间那些令人
振奋的讨论。

感谢克洛德·萨巴：感谢他的技术建议，感谢他把本书迎入自己主持的
丛书，尤其感谢他一直耐心地等候这本书。

\section*{英文版前言}

第一版（法文版）出版于2010年。此次新版：
\begin{itemize}
  \item 循Springer审稿人的建议，我们增补了几页关于\trMorseHomology{}与\trCellularHomology{}之关系的内容（即\cref{sec:cellular-homology}）。
  \item 循我们最初几位读者的建议，我们改正了数以百计（此处绝非夸张）
        的排印错误与疏失——其中，必须承认，颇有几个是数学性错误。%
        \footnote{这使我们在上文列举的数目之外又添了几项：本版现含
        七十四条定理、一百二十四条命题、一百一十一条引理与一百零二幅
        插图。}
\end{itemize}

我们对所有这些人深表感激，谨在此向他们致谢。

我们尤其感激 费利克斯·\trSchlenk，感谢他友好的建议与提问。还要再次
热忱感谢杜萨·\trMcDuff，感谢她善意的忠告与评注。

我们还更新了参考文献列表（参见第~11~章末关于\trLagrange{}\trSubmanifold{}的脚注），所增补的尤其是
本领域若干领军人物新近的书籍与论文：一篇关于\trConleyZehnderIndex
{}定义的非常清晰的论文 [43]、本书作者之一所获得的新结果 [23,~24]，
以及另一位作者关于\trArnoldConjecture{}历史的几点说明 [6]。

请注意，符号 $\vdash$ 既表示证明的结束，也表示未给出证明。%
\footnote{译注：原书此符号排作空心方框 $\square$；本译本统一以
$\vdash$ 复刻其职能。}\noqed

最后但同样重要的是：我们感谢雷妮·\trErne，她以非常的善意与耐心，
为一部相当技术性的文本完成了一项出色的翻译工作。

\begin{flushright}
  \trAudinFull\\
  \trDamianFull
\end{flushright}
